\subsection{SFERA}
Studien presenterar en IT-artefakt i form av förbättringsförslag som kunnat utrönas för SFERA-protokollet utifrån intervjuer med lokförare.
Förbättringsförslagen baserades på dokumentanalys av för ämnet relevanta dokument samt semistrukturerade intervjuer med fem lokförare.
Förbättringsförslagen valideras i en andra intervjuomgång av de fem respondenter som deltagit i den första intervjuomgången.
Observationer och forskning på praktisk tillämpning av SFERA saknas i publicerad vetenskaplig forskning i dag då protokollet, sannolikt då protokollet är publicerat relativt nyligen.
När SFERA-protokollet implementeras av järnvägsföretag
\subsection{Gamification och vidare forskning}
Vad gjordes i gamification-väg
De goda resultat från den något ringa forskning som bedrivits om gamification i digitala lokförarverktyg banar väg för mer djuplodande forskning i samarbete med branschen och speciellt järnvägsföretag som använder DAS- eller i synnerhet C-DAS-system.
Vikten av att lokförare följer sin körprofil är avgörande för att det digitala paradigmskiftet inom operativ järnvägsstyrning och lämpar sig väl för integration av gamificationprinciper.
Dock saknas praktiska tillämpningsexempel på gamification i C-DAS-sammanhang i forskningen.
Detta borde undersökas för att stödja järnvägsföretagen i deras tillämpning av gamification inom C-DAS.